\documentclass{ohm_project_description}

\projecttitle{Entwicklung eines mobilen Mini-Roboters}
\projectauthor{Labor für mobile Robotik}
\projectdate{2026}

\begin{document}

\maketitle
\thispagestyle{fancy}

Im Rahmen dieser Projektarbeit wird ein fahrbarer Mini-Roboter als Entwicklungs- und
Schulungsplattform konzipiert, aufgebaut und getestet. Ausgehend von einem bestehenden
Roboterkonzept entsteht eine eigenständige, mechanisch und elektrisch robuste Variante,
die kostengünstig nachgebaut werden kann – Materialkosten unter 1.000\,€ je Exemplar.
Der Roboter verfügt über einen omnidirektionalen Antrieb mit vier Mecanum-Rädern,
eine wiederaufladbare Energieversorgung mit mehreren Bordspannungen sowie eine
leistungsfähige Einplatinen-Steuerung mit aktiver Kühlung. Zur Umfelderkennung ist
ein 360°-LiDAR integriert; optional kann eine Tiefenkamera ergänzt werden. Die
Bedienung erfolgt wahlweise per Bluetooth-Controller oder installationsfreiem
Web-Interface vom Tablet.

\section*{Arbeitspakete}
\begin{itemize}[leftmargin=0.5cm]
    \setlength\itemsep{.1em}
    \item Konstruktion des mechanischen Aufbaus (mehrstöckige Aluminiumstruktur,
          Komponentenmontage mit niedrigem Schwerpunkt)
    \item Auswahl, Beschaffung und elektrische Integration der Komponenten
          (Antrieb, Energieversorgung mit DC-DC-Wandlung, Ladekonzept)
    \item Inbetriebnahme der Motorsteuerung und des omnidirektionalen Fahrantriebs
          (Mecanum-Räder, modulare Treiberplatinen)
    \item Integration und Test der Sensorik und Ausgabemedien (360°-LiDAR, Display,
          Lautsprecher, optional Tiefenkamera)
    \item Implementierung der Bedienschnittstellen (Bluetooth-Controller, Web-Interface,
          Datenspeicherung auf externem Datenträger)
    \item Dokumentation für den Nachbau sowie Kostenanalyse für die spätere
          Serienfertigung
\end{itemize}

\section*{Voraussetzungen}
\begin{itemize}[leftmargin=0.5cm]
    \setlength\itemsep{.1em}
    \item Interesse an mobiler Robotik und praktischem Systemaufbau
    \item Grundkenntnisse aus Elektrotechnik (Schaltungsaufbau, Energieversorgung)
    \item Grundkenntnisse aus Maschinenbau (mechanischer Aufbau, CAD von Vorteil)
    \item Grundkenntnisse in Programmierung (Python oder C++ hilfreich)
    \item Teamfähigkeit und Freude an interdisziplinärer Zusammenarbeit
    \item Geeignet für Studierende der Studiengänge BEEI, BEME und verwandter Richtungen
\end{itemize}

\vspace{0.5cm}
Das Thema kann nach Abstimmung als \textbf{Projektarbeit} (4 Studierende, ca.\ 4 Monate)
bearbeitet werden.

\vfill
\textcolor{ohm_red}{\rule{\linewidth}{0.4mm}}
\textbf{\textcolor{ohm_red}{Labor für mobile Robotik}} \\
\begin{tabular}{@{}ll}
\textbf{Betreuer:} & Prof.\ Dr.\ Christian Pfitzner \\
\textbf{E-Mail:}   & \href{mailto:christian.pfitzner@th-nuernberg.de}{christian.pfitzner@th-nuernberg.de} \\
\end{tabular}

\end{document}
